\documentclass[11pt]{article}
\usepackage{amsmath,amssymb,amsthm,mathtools}

\newtheorem{theorem}{Theorem}
\newtheorem{proposition}[theorem]{Proposition}
\newtheorem{lemma}[theorem]{Lemma}
\newtheorem{remark}[theorem]{Remark}

\DeclareMathOperator{\Pa}{Pa}

\begin{document}

\title{Scaling invariance of the angular lower-tail coefficient and an application to the Pareto-mixture copula}
\author{}
\date{}
\maketitle

\paragraph{Definition.}
Let $C$ be a copula on $[0,1]^2$. For $\delta\in(0,2)$ and $\gamma>0$, define
\[
\kappa_L(\delta,\gamma;C)
=
\lim_{u\downarrow 0}
\frac{\log C\!\bigl(\gamma u^\delta,\gamma^{-1}u^{2-\delta}\bigr)}{\log u},
\]
whenever the limit exists.

The purpose of this note is twofold:
\begin{enumerate}
    \item to prove that, under a minimal existence assumption, $\kappa_L(\delta,\gamma;C)$ does not depend on the scaling parameter $\gamma$;
    \item to verify that the Pareto-mixture copula $C_{\rm PM}$ satisfies this property.
\end{enumerate}

\section*{1. A general scaling-invariance theorem}

For the general statement below, no smoothness assumption on $C$ is needed beyond the defining properties of a copula. The only regularity condition is existence of the limit for one reference scaling.

\begin{theorem}\label{thm:gamma-invariance}
Fix $\delta\in(0,2)$ and let $C$ be a copula. Assume that there exists some $\tilde\gamma>0$ such that
\[
k
:=
\lim_{u\downarrow 0}
\frac{\log C\!\bigl(\tilde\gamma u^\delta,\tilde\gamma^{-1}u^{2-\delta}\bigr)}{\log u}
\]
exists as a finite real number. Then, for every $\gamma>0$,
\[
\lim_{u\downarrow 0}
\frac{\log C\!\bigl(\gamma u^\delta,\gamma^{-1}u^{2-\delta}\bigr)}{\log u}
=
k.
\]
Hence $\kappa_L(\delta,\gamma;C)$ is independent of $\gamma$.
\end{theorem}

\begin{proof}
Fix $\delta\in(0,2)$ and $\gamma>0$. Let
\[
r:=\frac{\gamma}{\tilde\gamma}>0,
\qquad
c_-:=\min\!\left\{r^{1/\delta},\,r^{-1/(2-\delta)}\right\},
\qquad
c_+:=\max\!\left\{r^{1/\delta},\,r^{-1/(2-\delta)}\right\}.
\]
Then $c_->0$ and $c_+>0$.

For every $u>0$,
\[
\frac{\tilde\gamma(c_-u)^\delta}{\gamma u^\delta}
=
\frac{\tilde\gamma c_-^\delta}{\gamma}
=
\frac{c_-^\delta}{r}
\le 1,
\]
because $c_-\le r^{1/\delta}$. Likewise,
\[
\frac{\tilde\gamma^{-1}(c_-u)^{2-\delta}}{\gamma^{-1}u^{2-\delta}}
=
\frac{\gamma}{\tilde\gamma}c_-^{2-\delta}
=
r\,c_-^{2-\delta}
\le 1,
\]
because $c_-\le r^{-1/(2-\delta)}$. Hence
\begin{equation}\label{eq:left-coord}
\tilde\gamma(c_-u)^\delta \le \gamma u^\delta,
\qquad
\tilde\gamma^{-1}(c_-u)^{2-\delta}\le \gamma^{-1}u^{2-\delta}.
\end{equation}

Similarly, since $c_+\ge r^{1/\delta}$ and $c_+\ge r^{-1/(2-\delta)}$,
\begin{equation}\label{eq:right-coord}
\tilde\gamma(c_+u)^\delta \ge \gamma u^\delta,
\qquad
\tilde\gamma^{-1}(c_+u)^{2-\delta}\ge \gamma^{-1}u^{2-\delta}.
\end{equation}

A copula is coordinatewise nondecreasing. Therefore \eqref{eq:left-coord} and \eqref{eq:right-coord} imply that, for every $u>0$,
\begin{equation}\label{eq:sandwich-C}
C\!\bigl(\tilde\gamma(c_-u)^\delta,\tilde\gamma^{-1}(c_-u)^{2-\delta}\bigr)
\le
C\!\bigl(\gamma u^\delta,\gamma^{-1}u^{2-\delta}\bigr)
\le
C\!\bigl(\tilde\gamma(c_+u)^\delta,\tilde\gamma^{-1}(c_+u)^{2-\delta}\bigr).
\end{equation}

Because the reference limit is assumed to exist as a finite real number, there exists $v_0>0$ such that
\[
C\!\bigl(\tilde\gamma v^\delta,\tilde\gamma^{-1}v^{2-\delta}\bigr)>0
\qquad\text{for all }0<v<v_0.
\]
Choose $u_0>0$ so small that $c_+u<v_0$ whenever $0<u<u_0$. Then all three terms in \eqref{eq:sandwich-C} are strictly positive for $u\in(0,u_0)$. Since $\log$ is increasing on $(0,\infty)$,
\[
\log C\!\bigl(\tilde\gamma(c_-u)^\delta,\tilde\gamma^{-1}(c_-u)^{2-\delta}\bigr)
\le
\log C\!\bigl(\gamma u^\delta,\gamma^{-1}u^{2-\delta}\bigr)
\le
\log C\!\bigl(\tilde\gamma(c_+u)^\delta,\tilde\gamma^{-1}(c_+u)^{2-\delta}\bigr).
\]
For $u$ sufficiently small, $\log u<0$. Dividing by $\log u$ reverses the inequalities:
\begin{equation}\label{eq:log-sandwich}
\frac{\log C\!\bigl(\tilde\gamma(c_-u)^\delta,\tilde\gamma^{-1}(c_-u)^{2-\delta}\bigr)}{\log u}
\ge
\frac{\log C\!\bigl(\gamma u^\delta,\gamma^{-1}u^{2-\delta}\bigr)}{\log u}
\ge
\frac{\log C\!\bigl(\tilde\gamma(c_+u)^\delta,\tilde\gamma^{-1}(c_+u)^{2-\delta}\bigr)}{\log u}.
\end{equation}

Now fix any $c>0$ and define
\[
A_c(u)
:=
\frac{\log C\!\bigl(\tilde\gamma(cu)^\delta,\tilde\gamma^{-1}(cu)^{2-\delta}\bigr)}{\log u}.
\]
For small $u>0$,
\[
A_c(u)
=
\frac{\log C\!\bigl(\tilde\gamma(cu)^\delta,\tilde\gamma^{-1}(cu)^{2-\delta}\bigr)}{\log(cu)}
\cdot
\frac{\log(cu)}{\log u}.
\]
As $u\downarrow 0$, we have $cu\downarrow 0$, so by the definition of $k$,
\[
\frac{\log C\!\bigl(\tilde\gamma(cu)^\delta,\tilde\gamma^{-1}(cu)^{2-\delta}\bigr)}{\log(cu)}
\longrightarrow k.
\]
Also,
\[
\frac{\log(cu)}{\log u}
=
1+\frac{\log c}{\log u}
\longrightarrow 1
\qquad(u\downarrow 0),
\]
because $\log u\to-\infty$. Therefore
\begin{equation}\label{eq:Ac-limit}
A_c(u)\longrightarrow k
\qquad(u\downarrow 0).
\end{equation}

Applying \eqref{eq:Ac-limit} with $c=c_-$ and $c=c_+$ shows that both outer terms in \eqref{eq:log-sandwich} converge to $k$. Hence, by the squeeze theorem,
\[
\lim_{u\downarrow 0}
\frac{\log C\!\bigl(\gamma u^\delta,\gamma^{-1}u^{2-\delta}\bigr)}{\log u}
=
k.
\]
This proves the claim.
\end{proof}

\begin{remark}
Theorem \ref{thm:gamma-invariance} is the precise formulation of the desired scaling invariance. For each fixed $\delta\in(0,2)$, it is enough to know that the limit exists for one value of $\gamma$; then it automatically exists for every $\gamma>0$ and has the same value.
\end{remark}

\section*{2. The Pareto-mixture copula}

We now verify the assumption of Theorem \ref{thm:gamma-invariance} for the proposed Pareto-mixture copula.

\paragraph{Model.}
Assume throughout that $\Pa(s)$ denotes the standard Pareto law on $[1,\infty)$ with survival function
\[
\mathbb P(Z>z)=z^{-s},\qquad z\ge 1.
\]
Let
\[
X_1=\frac{R}{R^*}\frac{H_1}{H_1^*},
\qquad
X_2=\frac{R}{R^*}\frac{H_2}{H_2^*},
\]
where
\[
R\sim\Pa(\beta),\quad R^*\sim\Pa(\alpha),\quad
H_1\sim\Pa(b_1),\quad H_1^*\sim\Pa(a_1),\quad
H_2\sim\Pa(b_2),\quad H_2^*\sim\Pa(a_2),
\]
all mutually independent, and all shape parameters are strictly positive.

Since $(X_1,X_2)$ has continuous margins, its copula is well defined; we denote it by $C_{\rm PM}$.

\subsection*{2.1. Logarithmic representation}

Set
\[
A:=\log R,\quad A^*:=\log R^*,\quad
B_i:=\log H_i,\quad B_i^*:=\log H_i^* \qquad (i=1,2).
\]
Then
\[
A\sim{\rm Exp}(\beta),\qquad A^*\sim{\rm Exp}(\alpha),
\qquad
B_i\sim{\rm Exp}(b_i),\qquad B_i^*\sim{\rm Exp}(a_i),
\]
independently, and
\[
-\log X_i = (A^*-A)+(B_i^*-B_i).
\]
Define
\[
U_0:=A^*-A,\qquad U_i:=B_i^*-B_i,\qquad i=1,2.
\]
Then
\[
-\log X_i = U_0+U_i,\qquad i=1,2.
\]

A direct convolution gives, for $s\in\mathbb R$,
\[
f_{U_0}(s)
=
\frac{\alpha\beta}{\alpha+\beta}
\begin{cases}
e^{-\alpha s}, & s\ge 0,\\[1mm]
e^{\beta s}, & s<0,
\end{cases}
\]
and for each $i=1,2$,
\[
\mathbb P(U_i>t)=\frac{b_i}{a_i+b_i}e^{-a_i t},
\qquad t\ge 0.
\]
Consequently there exist constants $0<k_i\le K_i<\infty$ such that, for all $z\in\mathbb R$,
\begin{equation}\label{eq:Ui-two-sided-bounds}
k_i e^{-a_i z_+}
\le
\mathbb P(U_i\ge z)
\le
K_i e^{-a_i z_+},
\qquad z_+:=\max\{z,0\}.
\end{equation}

\subsection*{2.2. Marginal lower-tail exponents}

Let $F_{X_i}$ denote the distribution function of $X_i$, and define
\[
m_i:=\min\{\alpha,a_i\},\qquad i=1,2.
\]

\begin{lemma}\label{lem:marginal-log-tail}
For each $i=1,2$,
\[
-\log F_{X_i}(e^{-t}) = m_i\,t + o(t)
\qquad (t\to\infty).
\]
Equivalently,
\[
F_{X_i}(e^{-t}) = \exp\!\bigl(-m_i t + o(t)\bigr).
\]
\end{lemma}

\begin{proof}
Fix $i\in\{1,2\}$. Since $X_i\le e^{-t}$ if and only if $U_0+U_i\ge t$,
\[
F_{X_i}(e^{-t})
=
\mathbb P(U_0+U_i\ge t)
=
\int_{\mathbb R}\mathbb P(U_i\ge t-s)\,f_{U_0}(s)\,ds.
\]
Split the integral into three parts:
\[
(-\infty,0],\qquad [0,t],\qquad [t,\infty).
\]
Using \eqref{eq:Ui-two-sided-bounds} and the formula for $f_{U_0}$, we obtain
\[
F_{X_i}(e^{-t})
\le
C_1 e^{-a_i t}
+
C_2 \int_0^t e^{-\alpha s}e^{-a_i(t-s)}\,ds
+
C_3 e^{-\alpha t},
\]
for suitable positive constants $C_1,C_2,C_3$ independent of $t$.

Hence
\[
F_{X_i}(e^{-t})
\le
C
\begin{cases}
e^{-a_i t}, & a_i<\alpha,\\[1mm]
(1+t)e^{-\alpha t}, & a_i=\alpha,\\[1mm]
e^{-\alpha t}, & \alpha<a_i,
\end{cases}
\]
for some constant $C>0$.

For the lower bound, we consider three cases.

\smallskip

\noindent
\emph{Case 1: $a_i<\alpha$.}
Integrating over $s\in[0,1]$ gives
\[
F_{X_i}(e^{-t})
\ge
c \int_0^1 e^{-\alpha s} e^{-a_i(t-s)}\,ds
\ge
c' e^{-a_i t}.
\]

\noindent
\emph{Case 2: $\alpha<a_i$.}
Integrating over $s\in[t,t+1]$ gives
\[
F_{X_i}(e^{-t})
\ge
c \int_t^{t+1} e^{-\alpha s}\,ds
\ge
c' e^{-\alpha t}.
\]

\noindent
\emph{Case 3: $\alpha=a_i$.}
Integrating over $s\in[0,t]$ gives
\[
F_{X_i}(e^{-t})
\ge
c \int_0^t e^{-\alpha s}e^{-\alpha(t-s)}\,ds
=
c\, t\, e^{-\alpha t}.
\]

Combining the upper and lower bounds yields
\[
F_{X_i}(e^{-t})
=
\exp\!\bigl(-m_i t + o(t)\bigr),
\]
which is the desired conclusion.
\end{proof}

\begin{lemma}\label{lem:quantile-log}
Let $F$ be continuous and strictly increasing near $0$, and suppose that for some $m>0$,
\[
-\log F(e^{-t}) = m t + o(t)
\qquad (t\to\infty).
\]
Then
\[
-\log F^{-1}(u)
=
\frac{1}{m}\log\!\frac{1}{u}
+
o\!\left(\log\!\frac{1}{u}\right)
\qquad (u\downarrow 0).
\]
\end{lemma}

\begin{proof}
Let $y(u):=-\log F^{-1}(u)$. Since $F(F^{-1}(u))=u$,
\[
-\log u
=
-\log F(e^{-y(u)})
=
m\,y(u)+o\!\bigl(y(u)\bigr).
\]
As $u\downarrow 0$, we have $F^{-1}(u)\downarrow 0$, hence $y(u)\to\infty$. Therefore
\[
y(u)
=
\frac{1}{m}\log\!\frac{1}{u}
+
o\!\left(\log\!\frac{1}{u}\right).
\]
\end{proof}

Applying Lemmas \ref{lem:marginal-log-tail} and \ref{lem:quantile-log} to $F_{X_1}$ and $F_{X_2}$, we obtain:
\begin{equation}\label{eq:quantile-asymptotics}
-\log F_{X_i}^{-1}(u)
=
\frac{1}{m_i}\log\!\frac{1}{u}
+
o\!\left(\log\!\frac{1}{u}\right),
\qquad i=1,2.
\end{equation}

\subsection*{2.3. Existence of the angular lower-tail limit for $C_{\rm PM}$}

Define
\[
q_1:=\frac{\delta}{m_1},
\qquad
q_2:=\frac{2-\delta}{m_2}.
\]

\begin{proposition}\label{prop:PM-existence}
For every $\delta\in(0,2)$, the limit
\[
\kappa_L(\delta,1;C_{\rm PM})
=
\lim_{u\downarrow 0}
\frac{\log C_{\rm PM}(u^\delta,u^{2-\delta})}{\log u}
\]
exists as a finite real number.

More precisely, let
\[
J(q_1,q_2)
:=
\inf_{s\ge 0}
\Bigl\{
\alpha s
+
a_1(q_1-s)_+
+
a_2(q_2-s)_+
\Bigr\}.
\]
Then
\[
\kappa_L(\delta,1;C_{\rm PM}) = J(q_1,q_2).
\]
\end{proposition}

\begin{proof}
Let
\[
u=e^{-t},\qquad t\to\infty.
\]
Define
\[
y_1(t):=-\log F_{X_1}^{-1}(u^\delta),
\qquad
y_2(t):=-\log F_{X_2}^{-1}(u^{2-\delta}).
\]
By \eqref{eq:quantile-asymptotics},
\begin{equation}\label{eq:yit}
y_1(t)=q_1 t+o(t),
\qquad
y_2(t)=q_2 t+o(t).
\end{equation}

Now
\[
C_{\rm PM}(u^\delta,u^{2-\delta})
=
\mathbb P\!\bigl(X_1\le F_{X_1}^{-1}(u^\delta),\,X_2\le F_{X_2}^{-1}(u^{2-\delta})\bigr).
\]
Since $X_i\le e^{-y_i(t)}$ is equivalent to $U_0+U_i\ge y_i(t)$, we obtain
\[
C_{\rm PM}(u^\delta,u^{2-\delta})
=
\mathbb P\!\bigl(U_0+U_1\ge y_1(t),\,U_0+U_2\ge y_2(t)\bigr).
\]
Conditioning on $U_0$ gives
\begin{equation}\label{eq:joint-cond}
C_{\rm PM}(u^\delta,u^{2-\delta})
=
\int_{\mathbb R}
\mathbb P\!\bigl(U_1\ge y_1(t)-s\bigr)\,
\mathbb P\!\bigl(U_2\ge y_2(t)-s\bigr)\,
f_{U_0}(s)\,ds.
\end{equation}

We first treat the case $q_1<q_2$. Then by \eqref{eq:yit}, for all sufficiently large $t$,
\[
y_1(t)\le y_2(t).
\]
Split the integral in \eqref{eq:joint-cond} into four intervals:
\[
(-\infty,0],\qquad [0,y_1(t)],\qquad [y_1(t),y_2(t)],\qquad [y_2(t),\infty).
\]
Using \eqref{eq:Ui-two-sided-bounds} and the explicit density of $U_0$, we obtain
\[
C_{\rm PM}(u^\delta,u^{2-\delta})
\le
C\Bigl(
e^{-A_t}+e^{-B_t}+e^{-D_t}
\Bigr),
\]
where
\[
A_t:=a_1y_1(t)+a_2y_2(t),
\qquad
B_t:=\alpha y_1(t)+a_2\bigl(y_2(t)-y_1(t)\bigr),
\qquad
D_t:=\alpha y_2(t),
\]
and $C>0$ is independent of $t$.

Indeed:
\begin{align*}
\int_{-\infty}^0 \cdots\,ds
&\le C e^{-a_1 y_1(t)-a_2 y_2(t)},\\
\int_0^{y_1(t)} \cdots\,ds
&\le C\!\left(e^{-a_1 y_1(t)-a_2 y_2(t)} + e^{-\alpha y_1(t)-a_2(y_2(t)-y_1(t))}\right),\\
\int_{y_1(t)}^{y_2(t)} \cdots\,ds
&\le C\!\left(e^{-\alpha y_1(t)-a_2(y_2(t)-y_1(t))} + e^{-\alpha y_2(t)}\right),\\
\int_{y_2(t)}^\infty \cdots\,ds
&\le C e^{-\alpha y_2(t)}.
\end{align*}

For the lower bound, we integrate over the unit intervals $[0,1]$, $[y_1(t),y_1(t)+1]$, and $[y_2(t),y_2(t)+1]$. Since the density $f_{U_0}$ is strictly positive on compact sets and \eqref{eq:Ui-two-sided-bounds} gives strictly positive lower bounds for the relevant tail probabilities, there exists $c>0$ such that
\[
C_{\rm PM}(u^\delta,u^{2-\delta})
\ge
c\,
\max\Bigl\{
e^{-A_t},\,e^{-B_t},\,e^{-D_t}
\Bigr\}.
\]
Combining upper and lower bounds,
\[
-\log C_{\rm PM}(u^\delta,u^{2-\delta})
=
\min\{A_t,B_t,D_t\}+O(1).
\]
Dividing by $t$ and using \eqref{eq:yit},
\[
-\frac{1}{t}\log C_{\rm PM}(u^\delta,u^{2-\delta})
\longrightarrow
\min\Bigl\{
a_1q_1+a_2q_2,\,
\alpha q_1+a_2(q_2-q_1),\,
\alpha q_2
\Bigr\}.
\]
A direct check shows that this minimum equals
\[
\inf_{s\ge 0}\Bigl\{\alpha s+a_1(q_1-s)_+ + a_2(q_2-s)_+\Bigr\}
=
J(q_1,q_2).
\]

If $q_2<q_1$, the same argument with the roles of indices $1$ and $2$ exchanged yields the same conclusion.

Finally, if $q_1=q_2$, then \eqref{eq:yit} implies
\[
y_1(t)-y_2(t)=o(t),
\]
so the two possible orderings differ only by $o(t)$ in the exponent. Hence the same limit still exists and equals $J(q_1,q_2)$.

Since $\log u=-t$, we conclude that
\[
\frac{\log C_{\rm PM}(u^\delta,u^{2-\delta})}{\log u}
=
-\frac{1}{t}\log C_{\rm PM}(u^\delta,u^{2-\delta})
\longrightarrow
J(q_1,q_2).
\]
Therefore
\[
\kappa_L(\delta,1;C_{\rm PM})=J(q_1,q_2),
\]
and the limit exists as a finite real number.
\end{proof}

\subsection*{2.4. Independence of the scaling parameter}

We can now combine Proposition \ref{prop:PM-existence} with Theorem \ref{thm:gamma-invariance}.

\begin{theorem}\label{thm:PM-gamma-independence}
For the Pareto-mixture copula $C_{\rm PM}$ and every $\delta\in(0,2)$,
\[
\kappa_L(\delta,\gamma;C_{\rm PM})
\]
exists for every $\gamma>0$ and is independent of $\gamma$.

In fact, if
\[
m_1=\min\{\alpha,a_1\},
\qquad
m_2=\min\{\alpha,a_2\},
\qquad
q_1=\frac{\delta}{m_1},
\qquad
q_2=\frac{2-\delta}{m_2},
\]
then for every $\gamma>0$,
\[
\kappa_L(\delta,\gamma;C_{\rm PM})
=
\inf_{s\ge 0}
\Bigl\{
\alpha s+a_1(q_1-s)_+ + a_2(q_2-s)_+
\Bigr\}.
\]
\end{theorem}

\begin{proof}
By Proposition \ref{prop:PM-existence}, for each fixed $\delta\in(0,2)$ the limit defining $\kappa_L(\delta,1;C_{\rm PM})$ exists as a finite real number. Therefore the hypothesis of Theorem \ref{thm:gamma-invariance} is satisfied with $\tilde\gamma=1$. Applying Theorem \ref{thm:gamma-invariance} yields
\[
\kappa_L(\delta,\gamma;C_{\rm PM})
=
\kappa_L(\delta,1;C_{\rm PM})
\qquad\text{for all }\gamma>0.
\]
The explicit formula then follows from Proposition \ref{prop:PM-existence}.
\end{proof}

\section*{Conclusion}

Under the minimal regularity condition that, for each fixed $\delta\in(0,2)$, the logarithmic limit exists for one reference scaling $\tilde\gamma>0$, the angular lower-tail coefficient
\[
\kappa_L(\delta,\gamma;C)
=
\lim_{u\downarrow 0}
\frac{\log C\!\bigl(\gamma u^\delta,\gamma^{-1}u^{2-\delta}\bigr)}{\log u}
\]
is independent of $\gamma$.

For the Pareto-mixture copula $C_{\rm PM}$, this condition is satisfied for every $\delta\in(0,2)$; hence $\kappa_L(\delta,\gamma;C_{\rm PM})$ exists for all $\gamma>0$ and is indeed scaling-invariant.

\end{document}